\scanpage{63}
\mspage{7}
% archive sheet number: 62
local'' $\overline F$ sur $S$, à valeurs dans $\mathrm{Hot}_W$, défini à
l'aide de $F$. Je rappelle qu'on trouve également
un tel syst. local. dans le cas d'une catégorie
Hot-fibrée sur $S$ (et $=$ un système local
à valeurs dans $\mathrm{Hot}$ \unsure{\}} lui-même, au lieu de
$\mathrm{Hot}_W$ \unsure{\}} -- on utilise ici tacitement que
$W_\infty\subset W$, i.e. \fbox{L3\,$|$\,quater}$^*$).
% slanted
\margin{On va le supposer vrai dans \unsure{toute|\deleted{tout}} la suite}
Et il en est de $=$ a fortiori pour les objets de
$\underline{\unsure{\mathrm{Parf}|\mathrm{Perf}}}(S)$. Pour $\mathrm{Liss}\,S$ resp. $\mathrm{Prop}.\,S$, on
trouve simplement un foncteur
\[
S^\circ\xrightarrow{\overline F}\mathrm{Hot}\qquad\text{resp.}\quad S\mathrel{\overset{\overline{F'}}{\unsure{-|\to}}}\mathrm{Hot},
\]
d'où en composant avec $\mathrm{Hot}\mathrel{-}\mathrm{Hot}_W$, des
foncteurs à valeurs dans $\mathrm{Hot}_W$.

Dans chacune de ces 14 catégories, on
introduit un ensemble de flèches, qui
serviront à localiser, et on le désignera tjrs
par la même notation $W_S$. Dans \deleted{toutes}
ces catégories sauf seulement $\mathrm{Cat}_{/S}$, \deleted{\unsure{la}}
$W_S$ désigne la $W$-équivalence ``fibre par fibre''.
Dans $\mathrm{Cat}_{/S}$ d'autre part, ce sera la
$W$-équivalence sans plus. toutes les flèches
du diagramme \unsure{p.|ci} sont compatibles avec les
ensembles localisants de flèches.
